\subsection{Criterios de exactitud de los funtores homológicos de módulos}
\label{subsection:III.7.3}

\begin{proposition}[7.3.1]
\label{III.7.3.1}
Sea $A$ un anillo, no necesariamente conmutativo, y sea
$T_\bullet$ un funtor homológico covariante
\emph{(T, II, 2.1)} de $\cat{Ab}_A$ en $\cat{Ab}$ que
conmuta con las sumas directas. Sea $p$ un entero tal que
$T_p$ y $T_{p-1}$ estén definidos. Las condiciones siguientes son
equivalentes:
\begin{enumerate}
  \item[(a)] $T_p$ es exacto por la derecha;
  \item[(b)] $T_{p-1}$ es exacto por la izquierda;

\oldpage[III]{181}

  \item[(c)] para todo $A$-módulo a izquierda $M$, el homomorfismo
    canónico y funtorial \eqref{III.7.2.2.1.eq}
    \[
      T_p(A_s)\otimes_A M
      \longrightarrow
      T_p(M)
      \tag{7.3.1.1}\label{III.7.3.1.1.eq}
    \]
    es un isomorfismo;
  \item[(d)] para todo $A$-módulo a izquierda $M$, el homomorfismo
    \eqref{III.7.3.1.1.eq} es un epimorfismo;
  \item[(e)] $T_p$ es isomorfo a un funtor de la forma
    \[
      M\longmapsto N\otimes_A M,
    \]
    donde $N$ es un $A$-módulo a derecha.
\end{enumerate}
Si, además, se verifican las condiciones de \sref{III.7.2.7} sobre
$A$ y $\mathfrak m$, las condiciones precedentes son también
equivalentes a
\begin{enumerate}
  \item[(f)] el homomorfismo canónico
    \[
      T_p(\varepsilon):T_p(A_s)\longrightarrow T_p(k)
    \]
    es un epimorfismo.
\end{enumerate}
\end{proposition}

\begin{proof}
Por definición de funtor homológico, $T_i$ es semiexacto para todo
$i$ para el cual esté definido; además, para toda sucesión exacta
\[
  0\longrightarrow M'
  \xrightarrow{u}M
  \xrightarrow{v}M''
  \longrightarrow0,
\]
se tiene
\[
  \Ker\bigl(T_{i-1}(u)\bigr)
  =
  \Coker\bigl(T_i(v)\bigr).
\]
Por tanto, es claro que (a) y (b) son equivalentes, y las demás
afirmaciones resultan inmediatamente de \sref{III.7.2.5} y
\sref{III.7.2.7}.
\end{proof}

\begin{corollary}[7.3.2]
\label{III.7.3.2}
Sea $A$ un anillo conmutativo. Con las notaciones de
\sref{III.7.3.1}, supongamos que $T_p$ es exacto por la derecha. Si
$f\in A$ no pertenece al anulador de ningún elemento no nulo de un
$A$-módulo $M$, entonces tampoco pertenece al anulador de ningún
elemento no nulo de $T_{p-1}(M)$. En particular, si $A$ es un
dominio entero, el $A$-módulo $T_{p-1}(A)$ es sin torsión.
\end{corollary}

\begin{proof}
Si $h_f$ designa la homotecia $x\mapsto fx$ de $M$, la hipótesis
significa que $h_f$ es inyectiva. La condición (b) de
\sref{III.7.3.1} muestra entonces que $T_{p-1}(h_f)$ también es
inyectiva.
\end{proof}

\begin{proposition}[7.3.3]
\label{III.7.3.3}
Sea $A$ un anillo, y sea $T_\bullet$ un funtor homológico
covariante de $\cat{Ab}_A$ en $\cat{Ab}$ que conmuta con las
sumas directas. Sea $p$ un entero tal que $T_{p-1}$, $T_p$ y
$T_{p+1}$ estén definidos. Las condiciones siguientes son
equivalentes:
\begin{enumerate}
  \item[(a)] $T_p$ es exacto;
  \item[(b)] $T_{p+1}$ y $T_p$ son exactos por la derecha;
  \item[(c)] $T_p$ y $T_{p-1}$ son exactos por la izquierda;
  \item[(d)] $T_{p+1}$ es exacto por la derecha y $T_{p-1}$ es
    exacto por la izquierda;
  \item[(e)] para todo $A$-módulo $M$, los homomorfismos canónicos
    \[
      T_i(A_s)\otimes_A M
      \longrightarrow
      T_i(M)
      \tag{7.3.3.1}\label{III.7.3.3.1.eq}
    \]
    son isomorfismos para $i=p$ e $i=p+1$;
  \item[($\mathrm{e}'$)] para todo $A$-módulo $M$, los
    homomorfismos \eqref{III.7.3.3.1.eq} son epimorfismos para
    $i=p$ e $i=p+1$;
  \item[(f)] para todo $A$-módulo $M$, el homomorfismo
    \eqref{III.7.3.3.1.eq} es un isomorfismo para $i=p$, y
    $T_p(A_s)$ es plano como $A$-módulo a derecha;
  \item[($\mathrm{f}'$)] para todo $A$-módulo $M$, el
    homomorfismo \eqref{III.7.3.3.1.eq} es un epimorfismo para
    $i=p$, y $T_p(A_s)$ es plano como $A$-módulo a derecha.
\end{enumerate}
\end{proposition}

\begin{proof}
La equivalencia de (a), (b), (c) y (d) resulta de la equivalencia
de las condiciones (a) y (b) de \sref{III.7.3.1}. La equivalencia
de (b), (e) y ($\mathrm{e}'$) resulta de la equivalencia de (a),
(c) y (d) en \sref{III.7.3.1}. Finalmente, decir que $T_p(A_s)$
es plano significa que el funtor
\[
  M\longmapsto T_p(A_s)\otimes_A M
\]
es exacto por la izquierda; la equivalencia de (a), (f) y
($\mathrm{f}'$) resulta de nuevo de la equivalencia de (a), (c) y
(d) en \sref{III.7.3.1}.
\end{proof}

\oldpage[III]{182}

\begin{corollary}[7.3.4]
\label{III.7.3.4}
Supongamos que $A$ es conmutativo, que $T_p$ es exacto y que
$T_p(A)$ es un $A$-módulo de presentación finita. Entonces la
función
\[
  x\longmapsto
  \operatorname{rank}_{\kres(x)}
  \bigl(T_p(\kres(x))\bigr)
\]
es localmente constante en $X=\Spec(A)$ y, por tanto, es constante
si $\Spec(A)$ es conexo.
\end{corollary}

\begin{proof}
Por la condición (f) de \sref{III.7.3.3}, $T_p(A)$ es un
$A$-módulo plano. Por consiguiente, es proyectivo de tipo finito, de
modo que $\widetilde{T_p(A)}$ es un $\sh O_X$-módulo localmente
libre (Bourbaki, \emph{Alg\`ebre commutative}, chap.~II, \S~5,
no.~2, th.~1). Además,
\[
  T_p(\kres(x))
  =
  T_p(A)\otimes_A\kres(x)
\]
por la condición (e) de \sref{III.7.3.3}, y se sabe que el rango en
$x$ del $A$-módulo $T_p(A)$ es localmente constante
\emph{(loc.\ cit.)}. Esto prueba el corolario.
\end{proof}

\begin{proposition}[7.3.5]
\label{III.7.3.5}
Supongamos que $A$ es un anillo artiniano a izquierda cuyo cociente
por su radical $\mathfrak m$ es un cuerpo $k$. Las condiciones de
\sref{III.7.3.3} son también equivalentes a cada una de las
siguientes:
\begin{enumerate}
  \item[(g)] el homomorfismo canónico
    \[
      T_i(\varepsilon):T_i(A_s)\longrightarrow T_i(k)
    \]
    es un epimorfismo para $i=p$ e $i=p+1$;
  \item[(h)] $T_p(\varepsilon)$ es un epimorfismo y
    $T_p(A_s)$ es plano como $A$-módulo a derecha o,
    equivalentemente, es un $A$-módulo libre a derecha (Bourbaki,
    \emph{Alg\`ebre commutative}, chap.~II, \S~3, no.~2,
    cor.~2 de la prop.~5).
\end{enumerate}
Supongamos además que $A$ es conmutativo y que el $A$-módulo
$T_p(k)$ tiene longitud finita $d$. Las condiciones precedentes son
entonces también equivalentes a cada una de las siguientes:
\begin{enumerate}
  \item[(i)] para todo $A$-módulo $M$ de longitud finita,
    \[
      \operatorname{length}\bigl(T_p(M)\bigr)
      =
      d\,\operatorname{length}(M);
      \tag{7.3.5.1}\label{III.7.3.5.1.eq}
    \]
  \item[(j)]
    \[
      \operatorname{length}\bigl(T_p(A)\bigr)
      =
      d\,\operatorname{length}(A).
      \tag{7.3.5.2}\label{III.7.3.5.2.eq}
    \]
\end{enumerate}
\end{proposition}

La equivalencia de (g) y (h) con las condiciones de
\sref{III.7.3.3} resulta inmediatamente de \sref{III.7.2.7}. Para
las demás afirmaciones utilizaremos el lema siguiente.

\begin{lemma}[7.3.5.3]
\label{III.7.3.5.3}
Sean $\cat{K}$ y $\cat{K}'$ dos categorías abelianas, y sea
\[
  F:\cat{K}\longrightarrow\cat{K}'
\]
un funtor covariante aditivo. Supongamos que $F$ es semiexacto y que
$F(S)$ tiene longitud finita en $\cat{K}'$ para todo objeto simple
$S$ de $\cat{K}$. Entonces $F(E)$ tiene longitud finita en
$\cat{K}'$ para todo objeto $E$ de longitud finita en $\cat{K}$.
Para toda sucesión exacta
\[
  0\longrightarrow E'
  \xrightarrow{u}E
  \xrightarrow{v}E''
  \longrightarrow0
\]
de objetos de longitud finita de $\cat{K}$, se tiene
\[
  \operatorname{length}F(E)
  \leq
  \operatorname{length}F(E')
  +
  \operatorname{length}F(E''),
  \tag{7.3.5.4}\label{III.7.3.5.4.eq}
\]
y se da la igualdad en \eqref{III.7.3.5.4.eq} si y solo si
\[
  0\longrightarrow F(E')
  \longrightarrow F(E)
  \longrightarrow F(E'')
  \longrightarrow0
\]
es exacta.
\end{lemma}

\begin{proof}
La sucesión
\[
  F(E')\xrightarrow{F(u)}F(E)\xrightarrow{F(v)}F(E'')
\]
es exacta por hipótesis. Si $F(E')$ y $F(E'')$ tienen longitud
finita, también la tienen $\Im(F(u))$ y $\Im(F(v))$. Puesto que
$\Ker(F(v))=\Im(F(u))$, el objeto $F(E)$ tiene longitud finita, y
\[
  \begin{split}
  \operatorname{length}F(E)
  &=
  \operatorname{length}\Im(F(u))
  +
  \operatorname{length}\Im(F(v))
  \\
  &\leq
  \operatorname{length}F(E')
  +
  \operatorname{length}F(E'').
  \end{split}
  \tag{7.3.5.5}\label{III.7.3.5.5.eq}
\]

\oldpage[III]{183}

Una inducción sobre la longitud de $E$ prueba la primera
afirmación. Además, la igualdad en \eqref{III.7.3.5.5.eq} solo
puede darse cuando
\[
  \operatorname{length}\Im(F(u))
  =
  \operatorname{length}F(E')
\]
y
\[
  \operatorname{length}\Im(F(v))
  =
  \operatorname{length}F(E'').
\]
La primera igualdad equivale a
$\operatorname{length}\Ker(F(u))=0$, y por tanto a $\Ker(F(u))=0$;
la segunda equivale a $\operatorname{length}\Coker(F(v))=0$, y por
tanto a $\Coker(F(v))=0$.
\end{proof}

Sea ahora $M$ un $A$-módulo de longitud finita, con $A$
conmutativo. Los cocientes de una serie de Jordan--Hölder de $M$
son necesariamente isomorfos al $A$-módulo $k$. Una inducción
sobre la longitud de $M$, usando \eqref{III.7.3.5.4.eq}, da por
tanto
\[
  \operatorname{length}T_p(M)
  \leq
  d\,\operatorname{length}(M).
  \tag{7.3.5.6}\label{III.7.3.5.6.eq}
\]
Se sigue también de \sref{III.7.3.5.3} que, si $T_p$ es exacto,
entonces se da la igualdad \eqref{III.7.3.5.1.eq}. Así, la condición
(a) de \sref{III.7.3.3} implica (i), y (i) implica claramente (j).
Queda por probar el lema siguiente.

\begin{lemma}[7.3.5.7]
\label{III.7.3.5.7}
La igualdad
\[
  \operatorname{length}T_p(A)
  =
  d\,\operatorname{length}A
\]
implica que $T_p(\varepsilon)$ es un epimorfismo y que $T_p(A)$ es
un $A$-módulo plano.
\end{lemma}

\begin{proof}
Partiendo de la sucesión exacta
\[
  0\longrightarrow\mathfrak m
  \longrightarrow A
  \longrightarrow k
  \longrightarrow0,
\]
las fórmulas \eqref{III.7.3.5.4.eq} y
\eqref{III.7.3.5.6.eq} dan
\[
  \begin{split}
  \operatorname{length}T_p(A)
  &\leq
  \operatorname{length}T_p(\mathfrak m)
  +
  \operatorname{length}T_p(k)
  \\
  &\leq
  d\bigl(
    \operatorname{length}\mathfrak m
    +
    \operatorname{length}k
  \bigr)
  =
  d\,\operatorname{length}A.
  \end{split}
\]
Por \sref{III.7.3.5.3}, la igualdad solo puede darse si
\[
  0\longrightarrow T_p(\mathfrak m)
  \longrightarrow T_p(A)
  \longrightarrow T_p(k)
  \longrightarrow0
  \tag{7.3.5.8}\label{III.7.3.5.8.eq}
\]
es exacta. Por \sref{III.7.2.7} y \sref{III.7.2.5}, $T_p$ es
isomorfo a un funtor
\[
  M\longmapsto N\otimes_A M.
\]
La exactitud de \eqref{III.7.3.5.8.eq}, junto con la sucesión exacta
de $\Tor$, muestra que
\[
  \Tor^A_1(N,k)=0.
\]
Se sigue que $N=T_p(A)$ es un $A$-módulo plano
\emph{(0, 10.1.3)}.
\end{proof}

\begin{lemma}[7.3.6]
\label{III.7.3.6}
Sea $A$ un anillo, y sea $T_\bullet$ un funtor homológico
covariante de $\cat{Ab}_A$ en $\cat{Ab}$ que conmuta con las
sumas directas. Supongamos que $T_p$ y $T_{p+1}$ están definidos y
que $T_p$ es exacto por la izquierda. Entonces $T_{p+1}$ es exacto
si y solo si $T_{p+1}(A_s)$ es plano como $A$-módulo a derecha.
\end{lemma}

\begin{proof}
Por \sref{III.7.3.1}, el homomorfismo canónico
\[
  T_{p+1}(A_s)\otimes_A M
  \longrightarrow
  T_{p+1}(M)
\]
es un isomorfismo de funtores. Basta ahora aplicar la definición de
$A$-módulo plano.
\end{proof}

\begin{proposition}[7.3.7]
\label{III.7.3.7}
Sea $A$ un anillo, y sea $T_\bullet$ un funtor homológico
covariante de $\cat{Ab}_A$ en $\cat{Ab}$ que conmuta con las
sumas directas. Supongamos que existe un $i_0$ tal que $T_i$ es
exacto para $i\leq i_0$. Entonces, para todo entero $p>i_0$, las
condiciones siguientes son equivalentes:
\begin{enumerate}
  \item[(a)] $T_q$ es exacto para $q\leq p$;
  \item[(b)] $T_q(A_s)$ es plano como $A$-módulo a derecha para
    $q\leq p$;
  \item[(c)] para todo $A$-módulo $M$, el homomorfismo canónico
    \[
      T_q(A_s)\otimes_A M
      \longrightarrow
      T_q(M)
    \]
    es sobreyectivo para $q\leq p+1$.
\end{enumerate}
\end{proposition}

\begin{proof}
La equivalencia de (a) y (b) resulta por inducción sobre $q$ de
\sref{III.7.3.6}, pues $T_{i_0}$ es exacto por hipótesis. La
equivalencia de (a) y (c) resulta de la equivalencia de las
condiciones (a) y ($\mathrm{e}'$) en \sref{III.7.3.3}.
\end{proof}

\oldpage[III]{184}

\begin{env}[7.3.8]
\label{III.7.3.8}
Sea $A$ un anillo conmutativo, sea $B$ una $A$-álgebra, no
necesariamente conmutativa, y sea $T_\bullet$ un funtor homológico
covariante de $\cat{Ab}_A$ en $\cat{Ab}$. Se sigue de las
definiciones de \sref{III.7.1.3} que el funtor de $\cat{Ab}_B$ en
$\cat{Ab}$ obtenido por extensión de escalares de $A$ a $B$,
denotado por
\[
  T_\bullet^{(B)}=\bigl(T_i^{(B)}\bigr),
\]
es también un funtor homológico.
\end{env}

\begin{corollary}[7.3.9]
\label{III.7.3.9}
Supongamos que $T_\bullet$ satisface las condiciones generales de
\sref{III.7.3.7} y conmuta con los límites inductivos. Supongamos
además que $A$ es un dominio entero y que cada $T_n(A)$ es un
$A$-módulo de presentación finita. Entonces, para todo entero $N$,
existe un $f\in A-\{0\}$ tal que
\[
  T_p^{(A_f)}:\cat{Ab}_{A_f}\longrightarrow\cat{Ab}
\]
es exacto para $p\leq N$.
\end{corollary}

\begin{proof}
Por hipótesis, $T_i$ es exacto para $i\leq i_0$, de modo que
$T_i(A)$ es plano para estos valores de $i$. Por la condición (b) de
\sref{III.7.3.7}, basta elegir $f$ de modo que
\[
  T_p^{(A_f)}(A_f)=T_p(A_f)
\]
sea un $A_f$-módulo libre para $i_0<p\leq N$. Ahora bien,
\[
  T_p(A_f)=\bigl(T_p(A)\bigr)_f,
\]
pues $T_p$ conmuta con los límites inductivos \sref{III.7.1.4}. Si
$x$ es el punto genérico de $\Spec(A)$, entonces
$\bigl(T_p(A)\bigr)_x$ es un espacio vectorial de dimensión finita
sobre el cuerpo de fracciones de $A$. Como cada $T_p(A)$ es de
presentación finita, existe un $f$ con la propiedad deseada
(Bourbaki, \emph{Alg\`ebre commutative}, chap.~II, \S~5,
no.~1, cor.\ de la prop.~2).

Si solo hay un número finito de índices $i$ tales que $T_i\neq0$,
se puede elegir $f\in A-\{0\}$ de modo que todos los
$T_p^{(A_f)}$ sean exactos.
\end{proof}

\begin{corollary}[7.3.10]
\label{III.7.3.10}
Supongamos que $T_\bullet$ satisface las condiciones generales de
\sref{III.7.3.7} y conmuta con los límites inductivos, que $A$ es
conmutativo y noetheriano, y que cada $T_n(A)$ es un $A$-módulo de
tipo finito. Entonces, para todo entero $N$, existe un abierto denso
$U$ de $\Spec(A)$ tal que, para todo entero $p\leq N$, la función
\[
  x\longmapsto
  \operatorname{rank}_{\kres(x)}
  \bigl(T_p(\kres(x))\bigr)
\]
es constante en $U$.
\end{corollary}

\begin{proof}
Sea $\mathfrak p$ un ideal primo minimal de $A$. El anillo
$B=A/\mathfrak p$ es un dominio entero, y $\Spec(B)$ se identifica
con una componente irreducible del espacio topológico $\Spec(A)$.
Probaremos por inducción sobre $p\leq N$ que existe un
$f_p\in B-\{0\}$ tal que, si $B'=B_{f_p}$, entonces
$T_i^{(B')}$ es exacto y $T_i(B')$ es un $B'$-módulo de tipo finito
para $i\leq p$.

La afirmación es cierta para $p\leq i_0$ por hipótesis, tomando
$f_p=1$, de modo que $B'=B=A/\mathfrak p$. En efecto, $T_p$ es
entonces exacto, y $T_p(B)$ es isomorfo a
\[
  T_p(A)/T_p(\mathfrak p);
\]
por tanto, es un $A$-módulo y, \emph{a fortiori}, un $B$-módulo de
tipo finito.

Para el paso de inducción, sea $f_p$ la imagen canónica en $B$ de un
elemento $g_p\in A$, y pongamos $A'=A_{g_p}$. Entonces
\[
  B'=A'/\mathfrak p',
\]
donde $\mathfrak p'$ es el ideal primo minimal
$\mathfrak pA_{g_p}$ de $A'$. Puesto que
\[
  T_i(A_{g_p})=\bigl(T_i(A)\bigr)_{g_p},
\]
los $T_i(A')$ son $A'$-módulos de tipo finito. Así,
$T_\bullet^{(A')}$ satisface las mismas hipótesis que $T_\bullet$,
con $i_0$ sustituido por $p$.

Podemos, por tanto, reducirnos al caso $A'=A$, en el que $T_p$ es
exacto. La sucesión exacta
\[
  0\longrightarrow\mathfrak p
  \longrightarrow A
  \longrightarrow A/\mathfrak p
  \longrightarrow0
\]
da entonces una sucesión exacta
\[
  T_{p+1}(A)
  \longrightarrow
  T_{p+1}(A/\mathfrak p)
  \xrightarrow{\partial}
  T_p(\mathfrak p)
  \longrightarrow
  T_p(A).
\]
Como $T_p$ es exacto, la última flecha es inyectiva. Por
consiguiente, $T_{p+1}(A/\mathfrak p)$ es un cociente de
$T_{p+1}(A)$ y es, por tanto, de tipo finito.

La demostración de \sref{III.7.3.9} solo ha utilizado la hipótesis
de tipo finito sobre $T_i(A)$ para $i\leq N$. Podemos, por tanto,
aplicar ese corolario al dominio entero $B$, al funtor
$T_\bullet^{(B)}$ y a $N=p+1$, lo que completa la inducción. Existe,
en consecuencia, un $f_N\in B-\{0\}$ tal que $\Spec(B_{f_N})$ es un
abierto $V$, denso en $\Spec(B)$, que no interseca ninguna otra
componente irreducible de $\Spec(A)$.

\oldpage[III]{185}

Si se ha probado la afirmación para $B_{f_N}$, existe un abierto
denso $W$ de $V$ en el que las funciones del enunciado son
constantes, pues
\[
  A_x=(B_{f_N})_x
  \qquad(x\in W).
\]
Repitiendo el mismo razonamiento para cada componente irreducible de
$\Spec(A)$ se prueba el corolario. Podemos, por tanto, reducirnos al
caso en que $A$ es un dominio entero. El razonamiento de
\sref{III.7.3.9} da entonces un $f\in A-\{0\}$ tal que los
$T_p(A_f)$ son $A_f$-módulos libres de tipo finito para $p\leq N$.
La conclusión resulta de \sref{III.7.3.4}.
\end{proof}

\begin{proposition}[7.3.11]
\label{III.7.3.11}
Sea $A$ un anillo local conmutativo de cuerpo residual $k$, y sea
$T_\bullet$ un funtor homológico covariante de $\cat{Ab}_A$ en
$\cat{Ab}$ que conmuta con las sumas directas. Supongamos que
existe un $i_0$ tal que $T_i$ es exacto para $i\leq i_0$, y que cada
$T_n(A)$ es un $A$-módulo de presentación finita. Las condiciones
equivalentes (a), (b) y (c) de \sref{III.7.3.7} implican las dos
condiciones siguientes, y son equivalentes a ellas cuando $A$ es
además reducido:
\begin{enumerate}
  \item[(d)] para todo $x\in\Spec(A)$,
    \[
      \operatorname{rank}_{\kres(x)}
      T_q(\kres(x))
      =
      \operatorname{rank}_k T_q(k)
      \qquad(q\leq p);
    \]
  \item[($\mathrm{d}'$)] para todo punto genérico $x_j$ de una
    componente irreducible de $\Spec(A)$,
    \[
      \operatorname{rank}_{\kres(x_j)}
      T_q(\kres(x_j))
      =
      \operatorname{rank}_k T_q(k)
      \qquad(q\leq p).
    \]
\end{enumerate}
\end{proposition}

\begin{proof}
Puesto que $T_q(A)$ es un $A$-módulo de presentación finita, la
condición (b) de \sref{III.7.3.7} equivale a decir que $T_q(A)$ es
libre para $q\leq p$ (Bourbaki, \emph{Alg\`ebre commutative},
chap.~II, \S~3, no.~2, cor.~2 de la prop.~5). La condición (c)
implica
\[
  T_q(\kres(x))
  =
  T_q(A)\otimes_A\kres(x)
  \qquad(q\leq p).
\]
Así, las condiciones equivalentes de \sref{III.7.3.7} implican (d),
y es inmediato que (d) implica ($\mathrm{d}'$).

Queda por mostrar que ($\mathrm{d}'$) implica (a) cuando $A$ es
reducido. Razonamos por inducción sobre $q\leq p$, pues $T_q$ es
exacto para $q\leq i_0$. Supongamos que $T_k$ es exacto para
$k\leq q<p$. Mostraremos que $T_{q+1}(A)$ es un $A$-módulo libre.
Por la hipótesis de inducción, el homomorfismo canónico
\[
  T_{q+1}(A)\otimes_A M
  \longrightarrow
  T_{q+1}(M)
\]
es un isomorfismo para todo $A$-módulo $M$, por la condición (c) de
\sref{III.7.3.7} y \sref{III.7.3.3}. Aplicando esto a
$M=\kres(x_j)$ y $M=k$, la condición ($\mathrm{d}'$) da
\[
  \operatorname{rank}_{\kres(x_j)}
  \bigl(T_{q+1}(A)\otimes_A\kres(x_j)\bigr)
  =
  \operatorname{rank}_k T_{q+1}(k)
\]
para todo $j$. Se sigue que $T_{q+1}(A)$ es libre (Bourbaki,
\emph{Alg\`ebre commutative}, chap.~II, \S~3, no.~2, prop.~7), lo
que completa la demostración.
\end{proof}

Los resultados precedentes se mejorarán considerablemente para los
funtores homológicos del tipo particular que estudiaremos en
\sref{III.7.4}: los criterios de exactitud así obtenidos solo hacen
intervenir uno de los $T_p$.
